\documentclass[aspectratio=169,11pt]{beamer}

\usetheme{default}
\usecolortheme{dove}
\usefonttheme{professionalfonts}

\setbeamertemplate{navigation symbols}{}
\setbeamertemplate{footline}[frame number]
\setbeamertemplate{itemize items}[circle]
\setbeamertemplate{enumerate items}[default]

\usepackage{amsmath,amssymb}
\usepackage{booktabs}
\usepackage{graphicx}
\usepackage{xcolor}

\definecolor{sbured}{RGB}{166,61,64}
\definecolor{sbugreen}{RGB}{31,122,91}
\definecolor{sbuamber}{RGB}{179,106,21}
\definecolor{sbublue}{RGB}{49,95,134}

\title{The Conditional Activation Model}
\subtitle{Uncovering Policy Demand Responses to AI-Driven Automation}
\author{Ignacio Urbina}
\institute{Stony Brook University}
\date{MPSA 2026 \(\cdot\) Job Market Paper}

\begin{document}

% ===================================================================
% TITLE
% ===================================================================
\begin{frame}
\titlepage
\end{frame}

% ===================================================================
% PUZZLE
% ===================================================================
\section{Puzzle}

% -------------------------------------------------------------------
\begin{frame}{Automation exposure is widespread, yet experimental evidence on policy demand remains inconclusive}

\begin{columns}[T]
\column{0.6\textwidth}
\begin{itemize}
    \item \textbf{Structural labor market shift} \\
    {\small Routine cognitive and manual tasks face increasing automation pressure}
    \item \textbf{Policy demand channel} \\
    {\small Standard models predict risk exposure $\to$ demand for retraining, regulation, redistribution}
    \item \textbf{Empirical gap} \\
    {\small Experimental tests of this channel produce inconsistent or null aggregate effects}
\end{itemize}

\column{0.35\textwidth}
\begin{block}{High-risk occupations}
\centering \textbf{\Large 47\%} \\
{\small of U.S.\ employment \\ (Frey \& Osborne 2017)}
\end{block}
\vspace{0.3cm}
\begin{block}{Public concern}
\centering \textbf{\Large 62\%} \\
{\small major concern about AI \\ workforce impact (Pew 2023)}
\end{block}
\end{columns}

\end{frame}

% -------------------------------------------------------------------
\begin{frame}{Three traditions converge: personal risk information should increase policy demand}

\begin{columns}[T]
\column{0.32\textwidth}
\begin{alertblock}{Self-interest}
\centering Risk $\to$ redistribution \\[4pt]
{\small Workers facing losses seek compensatory policy (Meltzer \& Richard 1981)}
\end{alertblock}

\column{0.32\textwidth}
\begin{alertblock}{Threat-motivated cognition}
\centering Anxiety $\to$ policy search \\[4pt]
{\small Threat perception focuses attention on protection (Albertson \& Gadarian 2015)}
\end{alertblock}

\column{0.32\textwidth}
\begin{alertblock}{Egotropic voting}
\centering Personal shock $\to$ preferences \\[4pt]
{\small Individual economic exposure drives policy demand (Kinder \& Kiewiet 1981)}
\end{alertblock}
\end{columns}

\end{frame}

% -------------------------------------------------------------------
\begin{frame}{Experimental results are inconsistent: null, positive, and opposing effects coexist}

\begin{itemize}
    \item \textbf{Null aggregate effects are common} \\
    {\small Zhang \& Hou (2022), Magistro et al.\ (2024): personalized risk information $\to$ null ATE on policy preferences}
    \item \textbf{Positive findings exist but with different designs} \\
    {\small Gallego et al.\ (2022): distinct treatment and sample yield positive effects}
    \item \textbf{Subgroup heterogeneity drives cancellation} \\
    {\small Arntz et al.\ (2024): opposing subgroup effects cancel in the aggregate. Wu (2022), Im (2023): effects moderated by ideology, education, prior beliefs}
\end{itemize}

\end{frame}

% ===================================================================
% THEORY
% ===================================================================
\section{Theory}

% -------------------------------------------------------------------
\begin{frame}{The unconditional ATE averages over heterogeneous subgroups}

\begin{columns}[c]
\column{0.45\textwidth}
\begin{block}{Standard estimand}
$E[Y(1) - Y(0)]$: unconditional ATE across full sample
\end{block}

\column{0.45\textwidth}
\begin{exampleblock}{Proposed estimand}
$E[Y(1) - Y(0) \mid S, \text{Gap}, \eta_{\text{anx}}]$: CATE stratified by trust type and activation state
\end{exampleblock}
\end{columns}

\vspace{0.5cm}
{\small Null aggregate effects may reflect composition rather than preference stability. The relevant estimand conditions on trust type and activation state.}

\end{frame}

% -------------------------------------------------------------------
\begin{frame}{CAM framework: two sequential gates}

\begin{columns}[T]
\column{0.45\textwidth}
\textbf{Processing sequence:}
\begin{enumerate}
    \item \textbf{Information} --- Expert risk estimate
    \item \textbf{Source trust} --- Accept or reject
    \item \textbf{Activation} --- Gap $\times$ Anxiety
    \item \textbf{Policy demand} --- Protection preferences
\end{enumerate}

\column{0.5\textwidth}
\begin{itemize}
    \item \textbf{Source trust: discrete classification} \\
    {\small Latent mixture separates Trusters from Skeptics --- not a continuous moderator}
    \item \textbf{Activation: joint requirement} \\
    {\small Neither belief gap nor trait anxiety alone is sufficient to trigger policy search}
    \item \textbf{Below-threshold pathway} \\
    {\small Beliefs update without policy preference change (cool cognition)}
\end{itemize}
\end{columns}

\end{frame}

% -------------------------------------------------------------------
\begin{frame}{Three qualitatively distinct cognitive pathways}

\begin{columns}[T]
\column{0.3\textwidth}
\begin{block}{\textcolor{sbured}{1. Source rejection}}
\centering \textbf{Skeptics} \\[4pt]
{\small Information fails the trust gate; no downstream processing regardless of gap or anxiety}
\end{block}

\column{0.3\textwidth}
\begin{block}{\textcolor{sbuamber}{2. Cool cognition}}
\centering \textbf{Below-threshold trusters} \\[4pt]
{\small Information is accepted and beliefs update, but insufficient vigilance to initiate policy search}
\end{block}

\column{0.3\textwidth}
\begin{block}{\textcolor{sbugreen}{3. Activated persuasion}}
\centering \textbf{Above-threshold trusters} \\[4pt]
{\small Both gates cleared; sustained elaboration of threat translates into demand for protective policy}
\end{block}
\end{columns}

\vspace{0.4cm}
\begin{center}
{\small \textit{If the three pathways are mixed in the population, the unconditional ATE will understate the responsive subgroup --- producing the null or weak aggregate effects observed in prior work.}}
\end{center}

\end{frame}

% ===================================================================
% DESIGN
% ===================================================================
\section{Design}

% -------------------------------------------------------------------
\begin{frame}{Pre-registered between-subjects experiment, $N = 901$ U.S.\ workers}

\begin{columns}[T]
\column{0.45\textwidth}
\textbf{Three-step belief correction protocol:}
\begin{enumerate}
    \item Explain expert methodology \\ {\small (comprehension: 93.4\%)}
    \item Elicit subjective risk belief \\ {\small (0--100 percentile)}
    \item Reveal expert estimate; display belief gap
\end{enumerate}

\vspace{0.3cm}
{\small \textcolor{sbugreen}{\textbf{Design property:}} Treatment intensity varies naturally via belief gap, creating continuous within-treatment variation.}

\column{0.5\textwidth}
\begin{itemize}
    \item \textbf{Outcome} \\
    {\small Policy demand index: 5 items, 7-pt scale ($\omega = 0.85$)}
    \item \textbf{Source trust} \\
    {\small 5-point item, post-treatment only}
    \item \textbf{Belief gap} \\
    {\small Expert percentile $-$ subjective prior}
    \item \textbf{Trait anxiety} \\
    {\small 4-item scale, pre-treatment ($\omega = 0.90$)}
\end{itemize}
\end{columns}

\end{frame}

% -------------------------------------------------------------------
\begin{frame}{Operationalization of the two-gate structure}

\begin{table}
\centering
\small
\begin{tabular}{lll}
\toprule
\textbf{Variable} & \textbf{Operationalization} & \textbf{Gate / Role} \\
\midrule
Policy demand ($Y$) & 5 items, 7-pt scale ($\omega = 0.85$) & Outcome \\
Source trust ($S$) & Single 5-pt item, post-treatment only & Gate 1: binary filter \\
Belief gap (Gap) & Expert percentile $-$ subjective prior & Gate 2: info.\ shock \\
Trait anxiety ($\eta$) & 4-item Big Five facet ($\omega = 0.90$) & Gate 2: vigilance \\
\bottomrule
\end{tabular}
\end{table}

\vspace{0.3cm}
{\small \textcolor{sbuamber}{\textbf{Asymmetry:}} Trust and gap are post-treatment (treated only); anxiety is pre-treatment (all respondents). The BFMM handles the resulting collider structure.}

\end{frame}

% -------------------------------------------------------------------
\begin{frame}{Manipulation check}

\begin{columns}[T]
\column{0.3\textwidth}
\begin{block}{Comprehension}
\centering \textbf{\Large 93.4\%} \\
{\small passed methodology check}
\end{block}

\column{0.3\textwidth}
\begin{block}{Mean belief gap}
\centering \textbf{\Large 2.44} \\
{\small SD = 36, range [$-89$, $+85$]}
\end{block}

\column{0.3\textwidth}
\begin{block}{Certainty $\times$ accuracy}
\centering \textbf{\Large $r = -0.004$} \\
{\small Metacognitive failure: confidence does not track accuracy}
\end{block}
\end{columns}

\end{frame}

% ===================================================================
% RESULTS
% ===================================================================
\section{Results}

% -------------------------------------------------------------------
\begin{frame}{Identification: observed trust is a post-treatment collider}

\begin{columns}[T]
\column{0.45\textwidth}
\textbf{Identification strategy:}
\begin{enumerate}
    \item \textbf{Problem} --- Gap $\to$ $Q_{\text{obs}}$ $\leftarrow$ $S$ creates collider bias
    \item \textbf{Block} --- Gap, Gap$^2$ block the collider path
    \item \textbf{Classify} --- Ordinal mixture $\to$ Trusters vs.\ Skeptics
\end{enumerate}

\vspace{0.3cm}
\begin{exampleblock}{Type-specific CATE}
\small
$\tau_T(\text{Gap}, \eta) = \beta_0^T + \beta_1^T \cdot \text{Gap} + \beta_2^T \cdot \eta + \beta_3^T \cdot \text{Gap} \times \eta$
\end{exampleblock}

\column{0.5\textwidth}
\textbf{BFMM architecture:}
\begin{itemize}
    \item \textbf{Trust layer} \\
    {\small Ordinal logistic mixture, Trusters vs.\ Skeptics, $\delta = 2.09$}
    \item \textbf{Measurement layer} \\
    {\small Anxiety CFA ($\omega = .90$), Policy CFA ($\omega = .85$)}
    \item \textbf{Activation layer} \\
    {\small Type-specific CATEs, Gap $\times$ $\eta_{\text{anx}}$ interaction, non-negativity}
\end{itemize}
\end{columns}

\end{frame}

% -------------------------------------------------------------------
\begin{frame}{Unconditional ATE is positive and robust ($d \approx 0.24$--$0.27$)}

\begin{table}
\centering
\small
\begin{tabular}{lccc}
\toprule
\textbf{Estimator} & \textbf{$d$} & \textbf{95\% CI} & \\
\midrule
OLS (unadjusted)    & 0.27 & [0.14, 0.40] & \textcolor{sbugreen}{$\bullet$} \\
OLS (adjusted)       & 0.24 & [0.13, 0.36] & \textcolor{sbugreen}{$\bullet$} \\
SEM                  & 0.24 & [0.13, 0.36] & \textcolor{sbugreen}{$\bullet$} \\
BFMM ($\pi$-weighted) & 0.27 & [0.14, 0.39] & \textcolor{sbugreen}{$\bullet$} \\
\bottomrule
\end{tabular}
\end{table}

\vspace{0.3cm}
{\small All specifications converge. The BFMM reproduces the ATE while decomposing it into type-specific components.}

\end{frame}

% -------------------------------------------------------------------
\begin{frame}{Source trust operates as a discrete gate}

\begin{table}
\centering
\small
\begin{tabular}{lcccc}
\toprule
 & \multicolumn{2}{c}{\textbf{Skeptics ($\pi \approx 62\%$)}} & \multicolumn{2}{c}{\textbf{Trusters ($\pi \approx 38\%$)}} \\
\cmidrule(lr){2-3} \cmidrule(lr){4-5}
 & Estimate & $P(\beta > 0)$ & Estimate & $P(\beta > 0)$ \\
\midrule
$\beta_0$: Intercept        & $+0.23$ & 0.98 & $+0.44$ & 1.00 \\
$\beta_1$: Gap               & $-0.06$ & 0.26 & $+0.14$ & 0.92 \\
$\beta_2$: Anxiety            & $-0.06$ & 0.23 & $+0.14$ & 0.93 \\
$\beta_3$: Gap $\times$ Anx   & $-0.07$ & 0.18 & $+0.13$ & 0.92 \\
\bottomrule
\end{tabular}
\end{table}

\vspace{0.3cm}
{\small \textcolor{sbuamber}{\textbf{Finding:}} Source rejection functions as a strict binary filter: all Skeptic moderation coefficients have $P(\beta > 0) < 0.30$; all Truster coefficients have $P(\beta > 0) > 0.92$.}

\end{frame}

% -------------------------------------------------------------------
\begin{frame}{Among Trusters, effects are jointly activated by Gap $\times$ Anxiety}

\begin{columns}[T]
\column{0.48\textwidth}
\begin{table}
\centering
\small
\begin{tabular}{lcc}
\toprule
\textbf{Condition} & \textbf{$d$} & \textbf{95\% CI} \\
\midrule
Unconditional ATE & 0.27 & [0.14, 0.39] \\
\midrule
High Gap $\times$ Low Anx  & 0.17 & [$-$0.13, 0.46] \\
Low Gap $\times$ High Anx  & 0.21 & [$-$0.07, 0.49] \\
Mean $\times$ Mean          & 0.29 & [0.18, 0.41] \\
\textcolor{sbugreen}{\textbf{High $\times$ High}} & \textcolor{sbugreen}{\textbf{0.61}} & \textcolor{sbugreen}{\textbf{[0.40, 0.86]}} \\
\bottomrule
\end{tabular}
\end{table}

\vspace{0.2cm}
{\footnotesize \textcolor{sbugreen}{\textbf{Pairwise contrasts (High $\times$ High vs.):}} \\
vs.\ High Gap $\times$ Low Anx: $P(\Delta > 0) = 0.98$ \\
vs.\ Low Gap $\times$ High Anx: $P(\Delta > 0) = 0.98$ \\
vs.\ Mean $\times$ Mean: $P(\Delta > 0) = 0.99$}

\column{0.48\textwidth}
\begin{itemize}
    \item \textbf{Gap alone insufficient} \\
    {\small High Gap $\times$ Low Anxiety: CI crosses zero ($d = 0.17$)}
    \item \textbf{Anxiety alone insufficient} \\
    {\small Low Gap $\times$ High Anxiety: CI crosses zero ($d = 0.21$)}
    \item \textbf{Joint activation required} \\
    {\small High $\times$ High: $d = 0.61$ --- 2.5$\times$ the unconditional ATE}
    \item \textbf{Activation threshold} \\
    {\small $\approx$ 20th percentile of trait anxiety ($z \approx -0.84$)}
\end{itemize}
\end{columns}

\end{frame}

% ===================================================================
% DISCUSSION
% ===================================================================
\section{Discussion}

% -------------------------------------------------------------------
\begin{frame}{Null aggregate effects reflect composition, not preference stability}

\begin{columns}[T]
\column{0.55\textwidth}
\begin{enumerate}
    \item \textbf{Theoretical:} Two-gate model resolves conflicting results as composition effects across three pathways.
    \item \textbf{Methodological:} BFMM addresses latent classification with collider-biased post-treatment moderators.
    \item \textbf{Empirical:} Maximum CATE ($d = 0.61$) exceeds ATE by factor 2.5; robust across specifications.
    \item \textbf{Implication:} Most workers dismiss expert warnings outright; among those who listen, policy demand activates only when the information surprises them and they are dispositionally vigilant.
\end{enumerate}

\column{0.4\textwidth}
\vspace{0.5cm}
{\small \textcolor{sbuamber}{\textbf{Scope:}} $N = 901$ U.S.\ workers $\cdot$ 62\% Democrat $\cdot$ immediate measurement $\cdot$ F\&O 2017 estimates. Generalizes to climate, vaccines, trade.}
\end{columns}

\end{frame}

\end{document}
